\section{System Architecture and Data Flow}

\subsection{Overall Architecture}

The Insight Financial Report Analysis Engine adopts a highly cohesive,
low-coupling modular design that perfectly adapts to the extreme local
computing power of Mac Studio (M3 Ultra). The system consists of four core
layers: data ingestion layer, retrieval and generation layer, evaluation
closed-loop layer, and interaction layer.

\subsection{System Architecture Diagram Description}

The system architecture diagram includes the following core components:

\begin{enumerate}
    \item \textbf{Data Ingestion Layer (Ingest)}: Responsible for PDF parsing, chunking, and vectorization
    \item \textbf{Retrieval \& Generation Layer}: Responsible for information retrieval and answer generation
    \item \textbf{Evaluation Layer}: Responsible for system performance evaluation
    \item \textbf{Interaction Layer (UI)}: Responsible for user interface and interaction
\end{enumerate}

Data flow: PDF report → Parsing → Chunking → Vectorization → Storage →
Retrieval → Reranking → Answer generation → Evaluation → Display

\subsection{Data Ingestion Layer (Ingest)}

The data ingestion layer is responsible for parsing, chunking, and vectorizing
financial PDF reports, providing basic data for subsequent retrieval and
generation.

\subsubsection{PDF Parsing}
PyMuPDFReader is used to parse PDF reports, extract text content, and preserve
page information. The system ensures that each document's metadata contains
`source` (file name) and `page\_label` (physical page number), which is the
foundation for implementing citation tracking.

\subsubsection{Chunking Strategy}
The system supports two chunking strategies:

\begin{enumerate}
    \item \textbf{Baseline Chunking (Phase 1)}: Uses a fixed-size chunking strategy with a chunk size of 256 tokens and 10% overlap.
    \item \textbf{Semantic Chunking (Phase 2)}: Based on BGE-M3 semantic chunking, which calculates cosine similarity between sentences and splits at semantic boundaries to maximize the preservation of contextual integrity for financial analysis.
\end{enumerate}

\subsubsection{Vectorization and Storage}
The BAAI/bge-m3 model is used to vectorize the chunked text, generate dense
vectors, and store them in a local ChromaDB vector database. For Phase 2, the
system also builds RAPTOR tree-based summaries, using background small models
to recursively cluster and summarize underlying semantic chunks to build a
macro knowledge tree.

\subsection{Retrieval \& Generation Layer}

The retrieval and generation layer is responsible for retrieving relevant
information from the vector database based on user queries and generating
answers through large language models.

\subsubsection{Retrieval Strategy}
The system supports two retrieval strategies:

\begin{enumerate}
    \item \textbf{Single-path Dense Vector Retrieval (Phase 1)}: Directly uses VectorStoreIndex for similarity retrieval.
    \item \textbf{Multi-path Hybrid Retrieval (Phase 2)}: Integrates BM25 retrieval and vector retrieval, fuses results through RRF (Reciprocal Rank Fusion), and uses bge-reranker-base for reranking.
\end{enumerate}

\subsubsection{Dialogue Engine}
The system uses CondensePlusContextChatEngine as the core dialogue engine,
integrating ChatMemoryBuffer to implement short-term memory management,
ensuring that the model can understand contextual references in multi-turn
dialogues.

\subsubsection{Generation Strategy}
The system uses locally deployed Qwen2.5-7B or Qwen2.5-72B models to generate
answers. During the generation process, the system requires the output of
citation information to ensure the verifiability of answers.

\subsection{Evaluation Layer}

The evaluation layer is responsible for evaluating the system's performance to
ensure the quality of system output.

The system uses the RAGAS evaluation framework, registering the locally running
Qwen2.5-72B as RAGAS's judge\_model to score system answers. Evaluation metrics
include:

\begin{enumerate}
    \item \textbf{Context Precision}: Measures whether document chunks that truly contain answer-related information are ranked at the top among the retrieved document chunks.
    \item \textbf{Faithfulness}: Measures whether the answer generated by the large language model is 100% derived from the retrieved document chunks without hallucinations.
    \item \textbf{Answer Relevancy}: Measures whether the answer directly addresses the user's original question.
\end{enumerate}

\subsection{Interaction Layer (UI)}

The interaction layer is responsible for building the user interface and
providing a friendly user experience.

The system uses Gradio to build the frontend interface, implementing the
following functions:

\begin{enumerate}
    \item Multi-document management: Supports uploading and managing multiple PDF
          reports.
    \item Conversational question-answering: Provides a natural language dialogue
          interface, supporting multi-turn conversations.
    \item Citation jump: Implements the function of clicking citations to jump to the
          corresponding page in the PDF.
    \item Artifact generation: Supports generating structured artifacts such as executive
          summaries and FAQs.
\end{enumerate}

\subsection{Data Flow}

The end-to-end data flow of the system is as follows:

\begin{enumerate}
    \item \textbf{Data Ingestion}: PDF report → PyMuPDF parsing → Chunking → BGE-M3 vectorization → ChromaDB storage
    \item \textbf{Retrieval \& Generation}: User query → Retrieval strategy distribution → Hybrid retrieval → Reranking → Large model answer generation
    \item \textbf{Evaluation}: System answer → RAGAS evaluation → Quantitative metric calculation
    \item \textbf{Interaction}: User input → Frontend processing → Backend processing → Frontend display
\end{enumerate}
